% Compile: pdflatex calculations.tex
\documentclass[11pt,a4paper]{article}
\usepackage[margin=1in]{geometry}
\usepackage{amsmath,amssymb,booktabs,xcolor,hyperref}
\hypersetup{colorlinks=true, linkcolor=black, urlcolor=blue!60!black}
\setlength{\parskip}{6pt}\setlength{\parindent}{0pt}

\title{\vspace{-2em}Trade Complementarity Index --- Calculations}
\author{}\date{}
\begin{document}\maketitle\vspace{-2em}

\section*{What this measures}
The TCI checks whether one country's exports match another country's imports. A high TCI means the two are natural trading partners. We compute it for each Indo-Pacific reporter against China and the US, first at HS6, then rolled up to HS4.

\section*{Notation}
HS chapters 84--90, years 2001--2024, all values in USD thousands.

\begin{center}\small
\begin{tabular}{ll}
\toprule
$X_{i,k}$ & Reporter $i$'s exports of product $k$ to the world \\
$X_{i}$   & Reporter $i$'s total exports \\
$X_{i,k}^{j}$ & Reporter $i$'s exports of $k$ to partner $j$ \\
$M_{j,k}$ & Partner $j$'s imports of $k$ from the world \\
$M_{j}$   & Partner $j$'s total imports \\
$W_{k}$   & World exports of $k$ \\
$W$       & World exports total \\
\bottomrule
\end{tabular}
\end{center}

\section*{Step 1 --- RCA at HS6}
RCA tells us whether a country sells more of product $k$ than the world average. Above 1 means specialized; below 1 means under-specialized.

\begin{equation*}
\mathrm{RCA}^{X}_{i,k} = \frac{X_{i,k}/X_{i}}{W_{k}/W}
\qquad
\mathrm{RCA}^{M}_{j,k} = \frac{M_{j,k}/M_{j}}{W_{k}/W}
\end{equation*}

Source: Balassa (1965).

\section*{Step 2 --- TCI at HS6}
TCI checks whether reporter $i$'s exports of $k$ line up with partner $j$'s imports of $k$. We compute it two ways and verify they match (algebraically they must --- if they don't, there's a bug).

\begin{equation*}
\mathrm{TCI}_{i,j,k} = \frac{X_{i,k}}{X_{i}} \cdot \frac{M_{j,k}}{M_{j}} \cdot \frac{W}{W_{k}}
\quad = \quad
\mathrm{RCA}^{X}_{i,k} \cdot \mathrm{RCA}^{M}_{j,k} \cdot \frac{W_{k}}{W}
\end{equation*}

\section*{Step 3 --- Roll up HS6 to HS4}

\textbf{Active filter.} An HS6 line counts only if both sides actually trade it: $X_{i,k}^{j} > 0$ AND $M_{j,k} > 0$. This keeps phantom codes from HS revisions out of the average.

\textbf{Weight.} Each active HS6 gets a weight equal to its share of the reporter's trade in that HS4 group:
\begin{equation*}
w_{i,j,k} = \frac{X_{i,k}}{\sum_{k' \text{ active in } h} X_{i,k'}}
\end{equation*}

\textbf{HS4 TCI.} Weighted average of the HS6 values:
\begin{equation*}
\mathrm{TCI}_{i,j,h} = \sum_{k \in h} w_{i,j,k} \cdot \mathrm{TCI}_{i,j,k}
\end{equation*}

We also compute TCI directly at HS4 from HS4-level totals as a second integrity check:
\begin{equation*}
\mathrm{TCI}^{RCA}_{i,j,h} = \mathrm{RCA}^{X}_{i,h} \cdot \mathrm{RCA}^{M}_{j,h} \cdot \frac{W_{h}}{W}
\end{equation*}

\section*{Open questions to resolve}
A few methodology and scope decisions are still open. The numbers in the current outputs reflect the choices noted below; if any of these change, the relevant tables need to be regenerated.

\begin{enumerate}\setlength{\itemsep}{2pt}
  \item \textbf{HS4 aggregation method.} The current rollup (Step 3) uses a weighted average with the reporter's HS6 export share as the weight. The strict Drysdale--Garnaut formulation sums the HS6 TCI values directly, with no second weighting. The two give different HS4 numbers (e.g., Korea--China 2020 HS 8542: 0.33 weighted average vs.\ 0.69 strict sum). Need to confirm which to use.
  \item \textbf{ICT scope.} The UNCTAD HS 2022 ICT goods classification covers 23 HS4 headings (107 HS6 codes). It is unclear whether to use the full UNCTAD list or a narrower subset focused on computers, semiconductors, and telecommunications.
  \item \textbf{IIT (Intra-Industry Trade) index.} A second index was requested alongside TCI but the formula has not been specified. The Grubel--Lloyd index is the standard choice; needs to be confirmed.
  \item \textbf{HS revision handling.} The dataset spans HS 1996, 2002, 2007, 2012, 2017, and 2022. Legacy codes (e.g., 8542.10 smart cards reclassified in HS 2007) carry residual values from late-adopting reporters past their nominal end dates. Currently retained at HS6; aggregated correctly at HS4. Confirm this is the intended treatment.
\end{enumerate}

\section*{Outputs}
For each partner (China, US), in \texttt{data/TradeMapData/export/}:
\begin{itemize}\setlength{\itemsep}{0pt}
  \item \texttt{\{partner\}\_TCI.xlsx} --- HS4 Summary (one row per Country/Year/HS4) and HS6 Detail (one row per Country/Year/HS6, with all inputs and weights).
  \item \texttt{\{partner\}\_\{HS4\}\_TCI\_RCA.png} --- TCI over time, one line per reporter.
\end{itemize}

\section*{Where the code lives}
All formulas are in \href{../loadFiles/services/TCICalculator.py}{\texttt{TCICalculator.py}}:

\begin{center}\small
\begin{tabular}{ll}
\toprule
\textbf{What} & \textbf{Function} \\
\midrule
Filter scope (chapters, years)   & \texttt{\_filter\_by\_hs\_chapter\_and\_year} \\
RCA at HS6                       & \texttt{\_calculate\_rca\_and\_tci\_rca} \\
TCI at HS6 (Drysdale--Garnaut form) & \texttt{\_calculate\_tci\_drysdale\_garnaut} \\
TCI at HS6 (RCA form)            & \texttt{\_calculate\_rca\_and\_tci\_rca} \\
HS4 weights and rollup           & \texttt{\_aggregate\_hs6\_to\_hs4} \\
Excel and chart output           & \texttt{\_export\_excel}, \texttt{\_export\_hs4\_tci\_charts} \\
\bottomrule
\end{tabular}
\end{center}

Input data is validated against UN Comtrade, and RCA values are validated against World Bank WITS. See \href{validation_methodology.md}{\texttt{validation\_methodology.md}}.

\end{document}